\chapter{Análisis de requisitos}
\label{chap:requisitos}

\ph{Este capítulo describe \textbf{qué} debe hacer el sistema, sin entrar todavía en cómo
lo hace. Es la base contractual del TFG: si un requisito no aparece aquí,
no debería darse por hecho que el sistema lo cumple.}

\section{Descripción del sistema}
\label{sec:descripcion_sistema}

\ph{Visión general del sistema en una página: qué es, qué resuelve y para quién.}

\ph{Cubre:}
\ph{
\begin{itemize}
  \item \ph{Tipo de sistema (aplicación web, CLI, librería, servicio\dots).}
  \item \ph{Qué entradas acepta y qué salidas produce.}
  \item \ph{Subsistemas o módulos principales en los que se divide.}
  \item \ph{Actores que interactúan con él (anónimo, autenticado, administrador\dots).}
\end{itemize}
}

\ph{Ejemplo de subsistemas:}
\ph{
\begin{itemize}
  \item \textbf{\ph{Subsistema 1}}: \ph{responsabilidad principal.}
  \item \textbf{\ph{Subsistema 2}}: \ph{responsabilidad principal.}
  \item \textbf{\ph{Subsistema 3}}: \ph{responsabilidad principal.}
\end{itemize}
}

\ph{Ejemplo de actores:}
\ph{
\begin{itemize}
  \item \textbf{\ph{Actor 1}}: \ph{rol y nivel de acceso.}
  \item \textbf{\ph{Actor 2}}: \ph{rol y nivel de acceso.}
\end{itemize}
}

El sistema desarrollado consiste en un generador automático de modelos de características integrado dentro del ecosistema de herramientas asociado a la variabilidad software. La solución consiste en una herramienta que sigue una arquitectura compatible con Flamapy y su integración dentro de UVLHub, permitiendo que los usuarios puedan configurar y ejecutar el proceso de generación mediante una interfaz web.

El objetivo principal del sistema es facilitar la creación automática de modelos de características expresados en UVL con distintos niveles de complejidad y expresividad. Estos modelos permiten disponer de conjuntos de prueba más completos para la evaluación, validación y mejora de herramientas relacionadas con la Ingeniería de Líneas de Productos Software (SPL), especialmente aquellas orientadas al análisis de variabilidad. De esta forma, el sistema aborda la necesidad de disponer de modelos con características avanzadas que permitan evaluar el comportamiento de dichas herramientas en diferentes escenarios.

La entrada principal del sistema consiste en una configuración definida por el usuario mediante los parámetros disponibles en la interfaz proporcionada por UVLHub. Estos parámetros permiten establecer las características del modelo que se desea generar, incluyendo aspectos relacionados con su estructura jerárquica, atributos y restricciones.

Como resultado, el sistema produce modelos expresados en el lenguaje UVL con los niveles de expresividad soportados. Estos modelos pueden visualizarse directamente tras finalizar el proceso de generación dentro de la plataforma UVLHub y también pueden descargarse como ficheros independientes con extensión \texttt{.uvl} para su posterior utilización en otras herramientas del ecosistema de variabilidad software.

El sistema puede dividirse principalmente en los siguientes subsistemas:

\begin{itemize}

    \item \textbf{Subsistema de configuración del generador:} responsable de gestionar los parámetros definidos por el usuario mediante la interfaz web proporcionada en UVLHub. Este subsistema permite establecer las características del modelo que se desea generar y preparar la configuración necesaria para ejecutar el proceso desde el navegador.

    \item \textbf{Subsistema de generación automática:} constituye el núcleo funcional del sistema y es el encargado de producir los modelos de características a partir de la configuración proporcionada, haciendo uso de los componentes y mecanismos disponibles dentro del framework Flamapy.

    \item \textbf{Subsistema de integración con UVLHub:} encargado de conectar la operación generadora desarrollada con el módulo \textit{generator} de UVLHub, permitiendo que la funcionalidad de la herramienta pueda ser utilizada dentro de la plataforma mediante los mecanismos definidos por esta.

\end{itemize}

Los actores principales que interactúan con el sistema son los siguientes:

\begin{itemize}

    \item \textbf{Usuario investigador o desarrollador:} configura los parámetros del generador y obtiene modelos UVL para su utilización en procesos de evaluación, experimentación o validación de herramientas de variabilidad software. El acceso al generador no requiere autenticación previa.

    \item \textbf{Herramientas de análisis de variabilidad:} utilizan los modelos generados como entrada para realizar diferentes tareas relacionadas con el análisis y procesamiento de modelos de características.

\end{itemize}

\section{Requisitos de información}
\label{sec:ri}

\ph{Describe las entidades de información que el sistema gestiona, con sus
atributos y restricciones (claves, unicidad, tipos, opcionalidad). Suele haber
una tabla por entidad.}

El sistema gestiona la información necesaria para configurar el proceso de generación automática de modelos de características. Esta información se representa mediante un conjunto de entidades de configuración que permiten definir los parámetros utilizados durante la generación.

La entidad principal del modelo conceptual es \textbf{FmgeneratorModel}, que agrupa las diferentes configuraciones específicas del generador. Cada una de estas configuraciones representa un aspecto concreto del modelo que se desea generar, como la nomenclatura, los niveles de expresividad UVL soportados, las características, la estructura jerárquica, las restricciones o los atributos.

Estas entidades son utilizadas posteriormente por la operación de generación automática para producir modelos de características expresados en UVL. La operación de generación no se considera una entidad de información, sino una funcionalidad del sistema que será descrita en los requisitos funcionales.

\subsubsection*{RI-01: FmgeneratorModel}

La entidad \textbf{FmgeneratorModel} representa la configuración global del generador automático de modelos de características. Contiene los parámetros generales del proceso y las configuraciones específicas asociadas a cada aspecto del modelo generado.

\begin{longtable}{P{4cm} P{3cm} P{6cm}}
\toprule
\textbf{Atributo} & \textbf{Tipo} & \textbf{Descripción} \\
\midrule
\endhead

num\_models & int & Número de modelos que se desean generar. \\

seed & int & Semilla utilizada para controlar la generación aleatoria. \\

ensure\_satisfiable & bool & Indica si los modelos generados deben ser satisfacibles. \\

naming & NamingConfig & Configuración relacionada con la nomenclatura de los elementos generados. \\

levels & LevelConfig & Configuración de los niveles de expresividad soportados por UVL. \\

features & FeaturesConfig & Configuración relacionada con la generación de características. \\

hierarchy & HierarchyConfig & Configuración de la estructura jerárquica del modelo. \\

constraints & ConstraintsConfig & Configuración de las restricciones generadas. \\

attributes & AttributesConfig & Configuración de los atributos asociados a las características. \\

\bottomrule
\caption{Atributos del requisito de información RI-01.}
\label{tab:ri-fmgenerator-model}
\end{longtable}


\subsubsection*{RI-02: NamingConfig}

La entidad \textbf{NamingConfig} representa la configuración utilizada para definir la nomenclatura de los elementos generados por el sistema.

\begin{longtable}{P{4cm} P{3cm} P{6cm}}
\toprule
\textbf{Atributo} & \textbf{Tipo} & \textbf{Descripción} \\
\midrule
\endhead

name\_prefix & str & Prefijo utilizado para los nombres generados. \\

include\_feature\_count\_suffix & bool & Indica si se incluye el número de características en el nombre del modelo generado. \\

include\_constraint\_count\_suffix & bool & Indica si se incluye el número de restricciones en el nombre del modelo generado. \\

\bottomrule
\caption{Atributos del requisito de información RI-02.}
\label{tab:ri-naming-config}
\end{longtable}


\subsubsection*{RI-03: LevelConfig}

La entidad \textbf{LevelConfig} representa los niveles de expresividad UVL que estarán disponibles durante el proceso de generación.

\begin{longtable}{P{4cm} P{3cm} P{6cm}}
\toprule
\textbf{Atributo} & \textbf{Tipo} & \textbf{Descripción} \\
\midrule
\endhead

boolean\_level & bool & Indica si se habilita la generación de modelos booleanos. \\

arithmetic\_level & bool & Indica si se habilitan elementos aritméticos en los modelos generados. \\

type\_level & bool & Indica si se permite el uso de tipos en características y atributos. \\

group\_cardinality & bool & Indica si se habilitan cardinalidades de grupo. \\

feature\_cardinality & bool & Indica si se habilitan cardinalidades individuales de características. \\

aggregate\_functions & bool & Indica si se permite la generación de funciones agregadas. \\

string\_constraints & bool & Indica si se habilitan restricciones sobre valores de tipo cadena. \\

\bottomrule
\caption{Atributos del requisito de información RI-03.}
\label{tab:ri-level-config}
\end{longtable}


\subsubsection*{RI-04: FeaturesConfig}

La entidad \textbf{FeaturesConfig} contiene los parámetros relacionados con la generación de características del modelo.

\begin{longtable}{P{4cm} P{3cm} P{6cm}}
\toprule
\textbf{Atributo} & \textbf{Tipo} & \textbf{Descripción} \\
\midrule
\endhead

min\_features & int & Número mínimo de características que puede contener el modelo generado. \\

max\_features & int & Número máximo de características que puede contener el modelo generado. \\

dist\_boolean & float & Probabilidad de generación de características booleanas. \\

dist\_integer & float & Probabilidad de generación de características enteras. \\

dist\_real & float & Probabilidad de generación de características reales. \\

dist\_string & float & Probabilidad de generación de características de tipo cadena. \\

min\_feature\_cardinality & int / list[int] & Valor mínimo permitido para las cardinalidades de características. \\

max\_feature\_cardinality & int / list[int] & Valor máximo permitido para las cardinalidades de características. \\

prob\_feature\_cardinality & float & Probabilidad de aplicar cardinalidades a las características. \\

\bottomrule
\caption{Atributos del requisito de información RI-04.}
\label{tab:ri-features-config}
\end{longtable}


\subsubsection*{RI-05: HierarchyConfig}

La entidad \textbf{HierarchyConfig} representa la configuración relacionada con la estructura jerárquica del modelo generado.

\begin{longtable}{P{4cm} P{3cm} P{6cm}}
\toprule
\textbf{Atributo} & \textbf{Tipo} & \textbf{Descripción} \\
\midrule
\endhead

max\_tree\_depth & int & Profundidad máxima permitida para la jerarquía del modelo. \\

dist\_optional & float & Probabilidad de generar relaciones opcionales. \\

dist\_mandatory & float & Probabilidad de generar relaciones obligatorias. \\

dist\_alternative & float & Probabilidad de generar relaciones alternativas. \\

dist\_or & float & Probabilidad de generar relaciones OR. \\

group\_cardinality\_min & int & Cardinalidad mínima para grupos de características. \\

group\_cardinality\_max & int & Cardinalidad máxima para grupos de características. \\

dist\_group\_cardinality & float & Probabilidad de generar relaciones con cardinalidad de grupo. \\

\bottomrule
\caption{Atributos del requisito de información RI-05.}
\label{tab:ri-hierarchy-config}
\end{longtable}


\subsubsection*{RI-06: ConstraintsConfig}

La entidad \textbf{ConstraintsConfig} contiene los parámetros utilizados para definir la generación de restricciones entre elementos del modelo.

\begin{longtable}{P{4cm} P{3cm} P{6cm}}
\toprule
\textbf{Atributo} & \textbf{Tipo} & \textbf{Descripción} \\
\midrule
\endhead

min\_constraints & int & Número mínimo de restricciones generadas. \\

max\_constraints & int & Número máximo de restricciones generadas. \\

extra\_constraint\_representativeness & int & Factor utilizado para aumentar la representatividad de las restricciones. \\

min\_vars\_per\_constraint & int & Número mínimo de variables utilizadas en una restricción. \\

max\_vars\_per\_constraint & int & Número máximo de variables utilizadas en una restricción. \\

ctc\_dist\_boolean & float & Probabilidad de generar restricciones booleanas. \\

ctc\_dist\_integer & float & Probabilidad de generar restricciones enteras. \\

ctc\_dist\_real & float & Probabilidad de generar restricciones reales. \\

ctc\_dist\_string & float & Probabilidad de generar restricciones sobre cadenas. \\

prob\_not & float & Probabilidad de aplicar negación en restricciones booleanas. \\

prob\_and & float & Probabilidad de utilizar operadores AND. \\

prob\_or\_ct & float & Probabilidad de utilizar operadores OR. \\

prob\_implication & float & Probabilidad de utilizar implicaciones. \\

prob\_equivalence & float & Probabilidad de utilizar equivalencias. \\

prob\_sum & float & Probabilidad de utilizar operaciones de suma. \\

prob\_substract & float & Probabilidad de utilizar operaciones de resta. \\

prob\_multiply & float & Probabilidad de utilizar operaciones de multiplicación. \\

prob\_divide & float & Probabilidad de utilizar operaciones de división. \\

prob\_equals & float & Probabilidad de utilizar comparaciones de igualdad. \\

prob\_less & float & Probabilidad de utilizar comparaciones menor que. \\

prob\_greater & float & Probabilidad de utilizar comparaciones mayor que. \\

prob\_less\_equals & float & Probabilidad de utilizar comparaciones menor o igual que. \\

prob\_greater\_equals & float & Probabilidad de utilizar comparaciones mayor o igual que. \\

prob\_sum\_function & float & Probabilidad de utilizar funciones agregadas de suma. \\

prob\_avg\_function & float & Probabilidad de utilizar funciones agregadas de promedio. \\

prob\_len\_function & float & Probabilidad de utilizar funciones de longitud sobre cadenas. \\

\bottomrule
\caption{Atributos del requisito de información RI-06.}
\label{tab:ri-constraints-config}
\end{longtable}


\subsubsection*{RI-07: AttributesConfig}

La entidad \textbf{AttributesConfig} representa la configuración relacionada con los atributos asociados a las características generadas.

\begin{longtable}{P{4cm} P{3cm} P{6cm}}
\toprule
\textbf{Atributo} & \textbf{Tipo} & \textbf{Descripción} \\
\midrule
\endhead

random\_attributes & bool & Indica si los atributos deben generarse automáticamente o definirse manualmente. \\

min\_attributes & int / None & Número mínimo de atributos generados en modo aleatorio. \\

max\_attributes & int / None & Número máximo de atributos generados en modo aleatorio. \\

\texttt{dist\_boolean\_atr} & float &
Probabilidad de seleccionar el tipo booleano en la generación
aleatoria de atributos. \\

\texttt{dist\_integer\_atr} & float &
Probabilidad de seleccionar el tipo entero en la generación
aleatoria de atributos. \\

\texttt{dist\_real\_atr} & float &
Probabilidad de seleccionar el tipo real en la generación
aleatoria de atributos. \\

\texttt{dist\_string\_atr} & float &
Probabilidad de seleccionar el tipo cadena en la generación
aleatoria de atributos. \\

attributes\_list & list[dict] &
Lista de atributos definidos manualmente para la generación. Cada
definición puede incluir su tipo, dominio, probabilidad de asociación y
participación en restricciones. \\

attribute\_attach\_probs & list[float] & Probabilidades de asociación de atributos a características. \\

attribute\_in\_constraints & list[bool] &
Indica si cada atributo definido manualmente puede participar en
restricciones. \\

\bottomrule
\caption{Atributos del requisito de información RI-07.}
\label{tab:ri-attributes-config}
\end{longtable}

\section{Requisitos funcionales}
\label{sec:rf}

\ph{Lista las capacidades concretas que el sistema debe ofrecer, agrupadas por
bloques temáticos (gestión de usuarios, lógica de negocio principal, exploración\dots).
Cada requisito tiene un identificador único, un nombre, una descripción
y una prioridad (Alta/Media/Baja).}

El sistema debe proporcionar las funcionalidades necesarias para permitir la configuración, ejecución y obtención de modelos de características generados automáticamente. Los requisitos funcionales identificados se muestran en la Tabla~\ref{tab:requisitos-funcionales}.

\begin{longtable}{P{1.5cm} P{4cm} P{6.5cm} P{2cm}}

\toprule
\textbf{ID} & \textbf{Nombre} & \textbf{Descripción} & \textbf{Prior.} \\
\midrule
\endhead

RF-01 & Configuración del generador &
El sistema debe permitir al usuario definir los parámetros necesarios para el proceso de generación, incluyendo aspectos relacionados con las características, niveles de expresividad, estructura jerárquica, restricciones y atributos del modelo generado.
& Alta \\[0.5em]

RF-02 & Validación de configuración &
El sistema debe comprobar que los parámetros introducidos por el usuario constituyen una configuración válida antes de iniciar el proceso de generación. Las restricciones y condiciones específicas aplicables a dicha validación se describen en la subsección \ref{sec:reglas_negocio}.
& Alta \\[0.5em]

RF-03 & Generación automática de modelos UVL &
El sistema debe permitir generar automáticamente modelos de características expresados en UVL a partir de la configuración proporcionada por el usuario.
& Alta \\[0.5em]

RF-04 & Soporte de niveles de expresividad UVL &
El sistema debe permitir generar modelos incorporando diferentes niveles de expresividad del lenguaje UVL, incluyendo construcciones booleanas, aritméticas y de tipos cuando hayan sido seleccionadas en la configuración.
& Alta \\[0.5em]

RF-05 & Visualización del modelo generado &
El sistema debe mostrar el modelo de características generado dentro de la plataforma UVLHub tras finalizar correctamente el proceso de generación.
& Media \\[0.5em]

RF-06 & Descarga del modelo generado &
El sistema debe permitir descargar el modelo generado como un fichero independiente con extensión \texttt{.uvl} para su posterior utilización en otras herramientas del ecosistema de variabilidad software.
& Alta \\

\bottomrule

\caption{Requisitos funcionales del sistema.}
\label{tab:requisitos-funcionales}

\end{longtable}

\section{Requisitos no funcionales}
\label{sec:rnf}

\ph{Atributos transversales del sistema: seguridad, rendimiento, mantenibilidad,
calidad de código, despliegue, estándares\dots Son tan vinculantes como los
funcionales y deben poder verificarse. Formula cada RNF con una MÉTRICA y un
UMBRAL comprobables, no con adjetivos vagos: en lugar de «el sistema debe ser
rápido», escribe «el 95\% de las peticiones se atiende en menos de 200~ms».
El Capítulo~\ref{chap:resultados} usará estos umbrales como criterios de
éxito medibles, así que un RNF que no se pueda medir aquí no se podrá
evaluar allí. El umbral de cada RNF cuantitativo debe coincidir con el del
criterio de éxito asociado de la Tabla~\ref{tab:criterios-exito}: un único
umbral, dos vistas.}


\begin{longtable}{P{1.5cm} P{4cm} P{7cm} P{2cm}}

\toprule
\textbf{ID} & \textbf{Nombre} & \textbf{Descripción} & \textbf{Prior.} \\
\midrule
\endhead

RNF-01 & Arquitectura modular y mantenibilidad &
El sistema debe mantener una estructura modular basada en la arquitectura de extensiones de Flamapy, siguiendo las convenciones establecidas por el ecosistema. Esta organización debe facilitar el mantenimiento del código, la incorporación de nuevas funcionalidades y la integración con otros componentes del framework.
& Alta \\[0.8em]

RNF-02 & Compatibilidad con modelos UVL &
Los modelos generados deben cumplir la especificación UVL soportada por las herramientas del ecosistema de variabilidad software, permitiendo su posterior utilización en procesos de análisis y evaluación de modelos de características.
& Alta \\[0.8em]

RNF-03 & Ejecución eficiente del proceso de generación &
El sistema debe minimizar la carga computacional sobre el servidor mediante la ejecución del proceso de generación en el entorno del cliente cuando sea posible. Para ello, se emplean tecnologías como Pyodide y WebAssembly (WASM), permitiendo ejecutar el código del generador directamente en el navegador y reduciendo los tiempos de respuesta asociados a la comunicación y procesamiento remoto.
& Alta \\[0.8em]

RNF-04 & Usabilidad de la interfaz de configuración &
La interfaz proporcionada para configurar el generador debe permitir al usuario seleccionar y comprender los parámetros disponibles de forma clara, facilitando la creación de configuraciones válidas y reduciendo posibles errores durante el proceso de configuración.
& Media \\[0.8em]

RNF-05 & Calidad del software &
El código desarrollado debe cumplir los mecanismos de control de calidad establecidos en el proyecto, incluyendo las comprobaciones automáticas mediante integración continua y las herramientas de análisis estático utilizadas para mejorar la legibilidad y mantenibilidad del código.
& Alta \\

\bottomrule

\caption{Requisitos no funcionales del sistema.}
\label{tab:requisitos-no-funcionales}

\end{longtable}

\section{Reglas de negocio}
\label{sec:reglas_negocio}

\ph{Las reglas de negocio (RN) son afirmaciones sobre cómo se comporta el
dominio que el sistema debe respetar, independientemente de la tecnología o
de la interfaz. No son requisitos funcionales (qué hace el sistema) ni
requisitos no funcionales (atributos transversales): son restricciones
lógicas del problema. Son las que se rompen últimas si cambian: no porque
hayamos refactorizado, sino porque el negocio ha cambiado.}

\begin{longtable}{P{1.5cm} P{4cm} P{9cm}}

\toprule
\textbf{ID} & \textbf{Nombre} & \textbf{Descripción} \\
\midrule
\endhead

RN-01 & Validación de parámetros generales &
Los parámetros generales del generador deben cumplir las restricciones mínimas necesarias para permitir la ejecución del proceso de generación. El número de modelos a generar (\texttt{num\_models}) debe encontrarse comprendido entre 1 y 1000. Además, la semilla utilizada para controlar la aleatoriedad del proceso (\texttt{seed}) debe ser un valor positivo.
\\[1em]

RN-02 & Validación de niveles de expresividad &
Los niveles de expresividad del generador deben respetar la jerarquía establecida entre niveles principales (\textit{major levels}) y niveles secundarios (\textit{minor levels}).

El nivel booleano (\texttt{boolean\_level}) constituye el nivel base del generador y debe permanecer siempre habilitado.

La activación de un nivel principal implica la disponibilidad de los niveles secundarios asociados. Por otro lado, un nivel secundario únicamente puede activarse cuando el nivel principal necesario para soportar su funcionalidad se encuentra habilitado.

En concreto, los niveles deben cumplir las siguientes dependencias:
\begin{itemize}
    \item El nivel aritmético (\texttt{arithmetic\_level}) permite la generación de expresiones numéricas y es necesario para habilitar los niveles secundarios relacionados con cardinalidades, funciones agregadas y operaciones aritméticas.
    \item El nivel de tipos (\texttt{type\_level}) requiere la activación previa del nivel aritmético.
    \item Los niveles secundarios no pueden permanecer activos si el nivel del que dependen está deshabilitado.
\end{itemize}
\\[1em]

RN-03 & Validación de configuración del árbol de características &
La configuración del árbol de características generado debe cumplir las restricciones necesarias para garantizar una estructura válida del modelo. Estas validaciones afectan al número de características, la estructura jerárquica y las distribuciones probabilísticas asociadas.
\\[0.5em]

& & Las restricciones establecidas son las siguientes:

\begin{itemize}[noitemsep, topsep=0pt]

    \item Los parámetros \texttt{min\_features} y \texttt{max\_features} deben ser mayores o iguales que uno, y el número mínimo de características no puede superar al número máximo configurado.

    \item La profundidad máxima del árbol (\texttt{max\_tree\_depth}) debe encontrarse comprendida entre uno y el número máximo de características configurado (\texttt{max\_features}).

    \item Las probabilidades asociadas a las relaciones jerárquicas deben encontrarse entre 0 y 1 y formar una distribución cuya suma sea exactamente uno.

    \item Cuando se habilite la cardinalidad de características (\texttt{feature\_cardinality}), los parámetros asociados deben cumplir las siguientes restricciones:
    \begin{itemize}[noitemsep, topsep=0pt]
        \item Los valores \texttt{min\_feature\_cardinality} y \texttt{max\_feature\_cardinality} deben ser mayores o iguales que uno, y el valor mínimo no puede superar al máximo.
        \item La probabilidad de asignación de cardinalidad (\texttt{prob\_feature\_cardinality}) debe encontrarse dentro del rango comprendido entre 0 y 1.
    \end{itemize}

\end{itemize}
\\[1em]

& &

\begin{itemize}[noitemsep, topsep=0pt]
    \item Cuando se habilite la cardinalidad de grupos (\texttt{group\_cardinality}), los parámetros asociados deben cumplir las siguientes restricciones:
    \begin{itemize}[noitemsep, topsep=0pt]
        \item Los valores \texttt{group\_cardinality\_min} y \texttt{group\_cardinality\_max} deben ser mayores o iguales que uno, y el valor mínimo no puede superar al máximo.
        \item Ambos valores no pueden superar el número máximo de características configurado (\texttt{max\_features}).
    \end{itemize}

    \item La distribución de tipos de características únicamente debe estar disponible cuando el nivel de tipos (\texttt{type\_level}) esté habilitado y sus probabilidades deben sumar exactamente uno.

\end{itemize}
\\[1em]

RN-04 & Validación de configuración de restricciones &
La configuración de restricciones entre características y atributos debe cumplir las siguientes reglas:

\begin{itemize}
    \item Los límites del número de restricciones (\texttt{min\_constraints} y \texttt{max\_constraints}) deben ser mayores o iguales que uno y el mínimo no puede superar al máximo.

    \item El número mínimo de variables utilizadas en una restricción (\texttt{min\_vars\_per\_constraint}) debe ser menor o igual que el número máximo permitido (\texttt{max\_vars\_per\_constraint}), siendo ambos valores mayores o iguales que uno.

    \item El parámetro \texttt{extra\_constraint\_representativeness} debe ser mayor o igual que uno.

    \item Las probabilidades utilizadas en la generación de restricciones deben encontrarse dentro del rango comprendido entre 0 y 1.

    \item La distribución asociada a los tipos de restricciones (\texttt{ctc\_dist\_boolean}, \texttt{ctc\_dist\_integer}, \texttt{ctc\_dist\_real} y \texttt{ctc\_dist\_string}) debe sumar exactamente uno cuando los niveles de expresividad correspondientes estén habilitados.

    \item La distribución de operadores booleanos (\texttt{prob\_and}, \texttt{prob\_or\_ct}, \texttt{prob\_implication} y \texttt{prob\_equivalence}) debe sumar exactamente uno. La probabilidad asociada al operador NOT (\texttt{prob\_not}) se considera de forma independiente y debe encontrarse dentro del rango comprendido entre 0 y 1.
\end{itemize}
\\[0.5em]

& &
\begin{itemize}
    \item Cuando el nivel aritmético esté habilitado y las funciones
    agregadas estén desactivadas, la distribución de operaciones
    aritméticas básicas (\texttt{prob\_sum}, \texttt{prob\_substract},
    \texttt{prob\_multiply} y \texttt{prob\_divide}) debe sumar uno.

    \item Cuando las funciones agregadas estén activadas, las probabilidades
    \texttt{prob\_sum\_function} y \texttt{prob\_avg\_function} se incorporan
    a la distribución de operaciones aritméticas. En este caso, la suma
    conjunta de las cuatro operaciones básicas y las dos funciones agregadas
    debe ser igual a uno, admitiendo una tolerancia absoluta de $10^{-6}$
    debida a la representación numérica.

    \item Cuando el nivel aritmético esté habilitado, la distribución de operadores de comparación (\texttt{prob\_equals}, \texttt{prob\_less}, \texttt{prob\_greater}, \texttt{prob\_less\_equals} y \texttt{prob\_greater\_equals}) debe sumar exactamente uno.
\end{itemize}
\\[1em]

RN-05 & Validación de configuración de atributos &
La configuración de atributos debe cumplir las restricciones asociadas al modo de generación seleccionado:

\begin{itemize}[noitemsep, topsep=0pt]

    \item En modo de generación automática (\texttt{random\_attributes}), los límites \texttt{min\_attributes} y \texttt{max\_attributes} deben estar definidos, deben ser mayores o iguales que cero y el mínimo no puede superar al máximo.

    \item Cuando se utiliza la generación automática de atributos,
    las probabilidades \texttt{dist\_boolean\_atr},
    \texttt{dist\_integer\_atr}, \texttt{dist\_real\_atr} y
    \texttt{dist\_string\_atr} deben encontrarse en el intervalo
    $[0,1]$ y formar una distribución cuya suma sea uno.

    \item En modo de definición manual, cada atributo debe cumplir las restricciones asociadas a su tipo:
    \begin{itemize}
        \item Los atributos booleanos deben definir al menos un valor posible.
        \item Los atributos de tipo entero, real o cadena deben establecer un valor mínimo y máximo válidos, siendo el mínimo menor o igual que el máximo.
    \end{itemize}

    \item En modo manual, los límites \texttt{min\_attributes} y
    \texttt{max\_attributes} deben tomar el valor \texttt{None}.

    \item En modo manual, la probabilidad de asociación de cada atributo,
    almacenada mediante el parámetro \texttt{attach\_probability}, debe
    pertenecer al intervalo $[0,1]$. Esta comprobación se realiza durante la
    validación de \texttt{AttributesConfig}.

\end{itemize}
\\[1em]

\bottomrule

\caption{Reglas de negocio del dominio.}
\label{tab:reglas-negocio}
\end{longtable}

\section{Restricciones del proyecto}
\label{sec:restricciones}

\ph{Limitaciones externas que acotan el espacio de soluciones del proyecto.
A diferencia de los RNF, no son atributos del sistema construido sino
condicionantes del proceso de construcción.}


\begin{longtable}{P{1.7cm} P{3.5cm} P{9.3cm}}

\toprule
\textbf{ID} & \textbf{Tipo} & \textbf{Descripción} \\
\midrule
\endhead

RES-01 & Tecnológica &
El desarrollo debe integrarse dentro del ecosistema existente de herramientas de variabilidad software, utilizando como base el framework Flamapy y la plataforma UVLHub. Esta condición limita la arquitectura de la solución, ya que el generador debe seguir los mecanismos de extensión definidos por Flamapy, adaptarse a la estructura proporcionada por UVLHub y utilizar las librerías propias del framework y componentes externos necesarios para la generación y representación de modelos de características.
\\[0.8em]

RES-02 & Tecnológica / interoperabilidad &
Los modelos generados deben utilizar el lenguaje UVL como formato de salida, debido a que constituye el estándar empleado dentro del ecosistema de herramientas con el que se integra el proyecto.
\\[0.8em]

RES-03 & Organizativa &
El desarrollo debe realizarse dentro del contexto del grupo de investigación y de los proyectos software existentes asociados, respetando los procedimientos, recursos y estructuras de trabajo establecidos.
\\[0.8em]

RES-04 & Temporal &
El proyecto está condicionado por la planificación académica del TFG, estableciendo un periodo limitado para realizar las fases de análisis, diseño, implementación, validación y documentación de la solución.
\\[0.8em]

\bottomrule

\caption{Restricciones del proyecto.}
\label{tab:res}
\end{longtable}

\section{Casos de uso}
\label{sec:casos_uso}

\subsection{Diagrama general de casos de uso}
\label{subsec:cu_diagrama}

\ph{Resumen narrativo del diagrama: qué actores hay, qué casos de uso conectan
con cada actor y qué relaciones \texttt{include} o \texttt{extend} existen
entre casos de uso.}

\ph{(Espacio reservado para el diagrama de casos de uso del sistema, no incluido
en esta plantilla.)}

\subsection{Diagrama general de casos de uso}
\label{subsec:cu_diagrama}

El diagrama de casos de uso representa las principales interacciones entre los actores externos y el sistema generador de modelos UVL. El usuario investigador o desarrollador es el encargado de configurar los parámetros del generador y, una vez finalizado el proceso, puede visualizar los modelos generados dentro de la plataforma o descargarlos para su posterior utilización. El sistema se encarga de validar la configuración proporcionada y ejecutar el proceso de generación automática de modelos de características UVL. Los modelos descargados pueden ser utilizados posteriormente por herramientas externas de análisis de variabilidad para realizar procesos de evaluación y análisis sobre los modelos generados.

La Figura~\ref{fig:diagrama-casos-uso} muestra el diagrama general de casos de uso del sistema.

\begin{figure}[H]
    \centering
    \includegraphics[width=\textwidth]{diagrama_casos_uso.png}
    \caption{Diagrama general de casos de uso del sistema. Elaboración propia}
    \label{fig:diagrama-casos-uso}
\end{figure}

\subsection{Especificación de casos de uso principales}
\label{subsec:cu_especificacion}

\ph{Especifica con detalle los casos de uso más relevantes. Cada uno usa una ficha
con los campos: actores, descripción, precondiciones, flujo principal,
flujos alternativos y postcondiciones. Repite la ficha tantas veces como casos
de uso quieras especificar.}


\begin{table}[H]
\centering
\small
\renewcommand{\arraystretch}{1.35}
\begin{tabular}{@{}>{\bfseries\RaggedRight\arraybackslash}p{4cm} P{10cm}@{}}
\toprule
\multicolumn{2}{l}{\textbf{CU-01 --- Configurar generador}} \\
\midrule

Actores &
Usuario investigador o desarrollador. \\

Descripción &
Permite al usuario definir la configuración necesaria para el proceso de generación automática de modelos de características UVL. El usuario puede seleccionar los niveles de expresividad y establecer los parámetros relacionados con características, estructura jerárquica, restricciones y atributos. \\

Precondiciones &
El usuario debe acceder a la interfaz de configuración del generador disponible dentro de UVLHub. \\

Flujo principal &
1. El usuario accede al formulario de configuración del generador.

2. El usuario establece los parámetros asociados a la generación del modelo en cada \textit{step}.

3. El sistema registra la configuración proporcionada. \\

Flujo alternativo &
Si los valores introducidos no cumplen las restricciones establecidas, el sistema informa al usuario de los errores detectados durante la introducción de dichos valores o antes de pasar al siguiente \textit{step}. \\

Postcondiciones &
Si la configuración de cada \textit{step} es válida, se procede en el último \textit{step} a la generación de los modelos.\\

\bottomrule
\end{tabular}
\caption{Especificación del caso de uso CU-01.}
\label{tab:uc01}
\end{table}


\begin{table}[H]
\centering
\small
\renewcommand{\arraystretch}{1.35}
\begin{tabular}{@{}>{\bfseries\RaggedRight\arraybackslash}p{4cm} P{10cm}@{}}
\toprule
\multicolumn{2}{l}{\textbf{CU-02 --- Validar configuración}} \\
\midrule

Actores &
Sistema generador de modelos UVL. \\

Descripción &
Permite comprobar que la configuración proporcionada por el usuario cumple las reglas de negocio definidas antes de ejecutar el proceso de generación automática. \\

Precondiciones &
Debe existir una configuración del generador proporcionada por el usuario. \\

Flujo principal &
1. El sistema recibe la configuración definida por el usuario y analiza sus parámetros.

3. El sistema comprueba las restricciones asociadas a características, jerarquía, restricciones y atributos en cada \textit{step}, respectivamente.

4. El sistema valida las distribuciones probabilísticas y dependencias entre parámetros. \\

Flujo alternativo &
Si alguna regla de negocio no se cumple, el sistema devuelve los errores detectados e impide continuar con la generación hasta corregir la configuración. \\

Postcondiciones &
Si la configuración de cada \textit{step} es válida, se le permite al usuario pasar al siguiente \textit{step}. \\

\bottomrule
\end{tabular}
\caption{Especificación del caso de uso CU-02.}
\label{tab:uc02}
\end{table}


\begin{table}[H]
\centering
\small
\renewcommand{\arraystretch}{1.35}
\begin{tabular}{@{}>{\bfseries\RaggedRight\arraybackslash}p{4cm} P{10cm}@{}}
\toprule
\multicolumn{2}{l}{\textbf{CU-03 --- Generar modelo de características UVL}} \\
\midrule

Actores &
Sistema generador de modelos UVL. \\

Descripción &
Permite generar automáticamente un modelo de características expresado en UVL a partir de una configuración previamente validada. \\

Precondiciones &
Las dependencias necesarias deben ser cargadas correctamente al acceder al primer \textit{step}. Además, la configuración proporcionada por el usuario debe ser validada en cada \textit{step}. \\

Flujo principal &
1. Cuando el usuario llega al último \textit{step}, el sistema inicializa el proceso de generación utilizando la configuración proporcionada.

2. El sistema genera, para cada modelo y de forma aleatoria, la estructura jerárquica del modelo de características, incorporando los niveles de expresividad, atributos y restricciones pertinentes, cumpliendo con la configuración seleccionada por el usuario.

3. El sistema genera la representación final de los modelos en formato UVL y los almacena temporalmente para su visualización o descarga. \\

Flujo alternativo &
Si se produce un error durante el proceso de generación, el sistema informa del fallo y no genera los modelos solicitados. \\

Postcondiciones &
Existe un/unos modelo/s de características UVL generado/s correctamente disponible/s para su posterior visualización o descarga. \\

\bottomrule
\end{tabular}
\caption{Especificación del caso de uso CU-03.}
\label{tab:uc03}
\end{table}


\begin{table}[H]
\centering
\small
\renewcommand{\arraystretch}{1.35}
\begin{tabular}{@{}>{\bfseries\RaggedRight\arraybackslash}p{4cm} P{10cm}@{}}
\toprule
\multicolumn{2}{l}{\textbf{CU-04 --- Visualizar modelo generado}} \\
\midrule

Actores &
Usuario investigador o desarrollador. \\

Descripción &
Permite al usuario consultar visualmente el modelo de características generado dentro de la plataforma UVLHub. \\

Precondiciones &
Debe existir al menos un modelo generado correctamente por el sistema. \\

Flujo principal &
1. El sistema carga la representación del modelo generado.

2. El sistema muestra el modelo dentro de la interfaz de UVLHub al hacer clic en el botón correspondiente.

3. El usuario inspecciona la estructura y elementos generados. \\

Flujo alternativo &
Si no se ha podido realizar la generación previa a la visualización, el sistema lanza un error informativo. \\

Postcondiciones &
El usuario puede consultar el modelo generado dentro de la plataforma. \\

\bottomrule
\end{tabular}
\caption{Especificación del caso de uso CU-04.}
\label{tab:uc04}
\end{table}


\begin{table}[H]
\centering
\small
\renewcommand{\arraystretch}{1.35}
\begin{tabular}{@{}>{\bfseries\RaggedRight\arraybackslash}p{4cm} P{10cm}@{}}
\toprule
\multicolumn{2}{l}{\textbf{CU-05 --- Descargar modelo generado}} \\
\midrule

Actores &
Usuario investigador o desarrollador. \\

Descripción &
Permite descargar en formato zip los modelos generados como ficheros independientes en formato UVL para su posterior utilización dentro o fuera de la plataforma. \\

Precondiciones &
Debe existir un modelo de características generado correctamente. \\

Flujo principal &
1. El usuario solicita la descarga del modelo generado haciendo clic en el botón correspondiente.

2. El sistema prepara el zip con los ficheros UVL generados. El usuario puede ver el archivo en la bandeja de descargas de su navegador.
\\

Flujo alternativo &
Si el modelo no se encuentra disponible o se produce un error durante la generación del fichero, el sistema informa al usuario y cancela la descarga. \\

Postcondiciones &
El modelo UVL queda disponible como fichero independiente para su almacenamiento o análisis mediante herramientas externas. \\

\bottomrule
\end{tabular}
\caption{Especificación del caso de uso CU-05.}
\label{tab:uc05}
\end{table}

\section{Matriz de trazabilidad}
\label{sec:trazabilidad}

\ph{La matriz de trazabilidad cruza los requisitos funcionales con los
casos de uso que los implementan y, opcionalmente, con las pruebas que los
verifican. Demuestra que ningún requisito queda huérfano y que cada caso
de uso responde a uno o más requisitos identificados.}

\ph{Marca con \texttimes\ la celda en la que el caso de uso (o la prueba)
cubre el requisito de la fila. Una fila vacía es un requisito sin cobertura;
una columna vacía es un caso de uso sin requisito que lo justifique.
Ambos son señales de alerta.}

\subsection{Requisitos funcionales \texorpdfstring{$\leftrightarrow$}{<->} casos de uso}
\label{subsec:trazabilidad_rf_cu}

\ph{Ajusta el número de columnas al número real de casos de uso
especificados (la tabla de ejemplo muestra seis columnas, pero solo se han
incluido tres fichas de ejemplo).}

\begin{longtable}{P{1.7cm} *{5}{P{1.8cm}}}

\toprule
\textbf{} & \textbf{CU-01} & \textbf{CU-02} & \textbf{CU-03} & \textbf{CU-04} & \textbf{CU-05} \\
\midrule
\endhead

RF-01 & \texttimes & & & & \\

RF-02 & & \texttimes & & & \\

RF-03 & & & \texttimes & & \\

RF-04 & & & \texttimes & & \\

RF-05 & & & & \texttimes & \\

RF-06 & & & & & \texttimes \\

\bottomrule

\caption{Matriz de trazabilidad entre requisitos funcionales y casos de uso.}
\label{tab:trazabilidad_rf_cu}

\end{longtable}

\subsection{Requisitos \texorpdfstring{$\leftrightarrow$}{<->} pruebas}
\label{subsec:trazabilidad_rf_pruebas}

\ph{Cruce de los requisitos con la suite de pruebas (Capítulo~\ref{chap:pruebas}).
Cada requisito debería tener al menos una prueba que lo cubra. Indica el
identificador o nombre del test correspondiente.}


\begin{longtable}{P{1.5cm} P{3.2cm} P{9.3cm}}

\toprule
\textbf{Requisito} & \textbf{Tipo de verificación} & \textbf{Evidencia} \\
\midrule
\endhead

RF-01 &
Integración / sistema &
\texttt{test\_full\_happy\_path\_produces\_valid\_params}
(\texttt{test\_wizard\_flow.py}), que comprueba que el usuario puede
completar el flujo de configuración y establecer los parámetros necesarios
para la generación. \\[0.8em]

RF-02 &
Pruebas unitarias &
\texttt{test\_step1\_rejects\_batch\_size\_that\_could\_break\_generation},
\texttt{test\_step2\_rejects\_minor\_levels\_without\_required\_major\_levels}
y
\texttt{test\_step3\_rejects\_invalid\_relation\_distribution}
(\texttt{test\_validators.py}), que comprueban el rechazo de configuraciones
inválidas. \\[0.8em]

RF-03 &
Prueba de extremo a extremo &
\texttt{test\_step6\_generates\_one\_uvl\_model\_with\_pyodide}
(\texttt{test\_selenium.py}), que verifica que el proceso de generación se
completa desde la interfaz y produce un modelo. \\[0.8em]

RF-04 &
Pruebas de integración &
\texttt{test\_everything\_on\_every\_family\_represented}
(\texttt{test\_wizard\_outputs.py}), que comprueba que las configuraciones
relacionadas con los diferentes niveles de expresividad se reflejan en los
modelos generados. \\[0.8em]

RF-05 &
Prueba de extremo a extremo &
\texttt{test\_step6\_generates\_one\_uvl\_model\_with\_pyodide}
(\texttt{test\_selenium.py}), considerando en este caso la comprobación de
que el modelo generado se muestra correctamente en la interfaz. \\[0.8em]

RF-06 &
Pruebas de integración &
\texttt{test\_filename\_suffixes\_applied\_to\_generated\_files}
(\texttt{test\_wizard\_outputs.py}), que comprueba la aplicación de los
sufijos configurados a los ficheros generados. \\

\midrule

RNF-01 &
Inspección estructural e integración &
Revisión de la separación entre el modelo de configuración, las clases
especializadas, la operación \texttt{GenerateFeatureModel} y la integración
con Flamapy y UVLHub. Esta revisión se complementa con la ejecución de las
pruebas de integración del módulo \texttt{generator}. \\[0.8em]

RNF-02 &
Validación de formato y prueba de sistema &
La prueba
\texttt{test\_step6\_generates\_one\_uvl\_model\_with\_pyodide}
(\texttt{test\_selenium.py}) comprueba que el resultado se genera y se
entrega en formato UVL. La validez sintáctica del modelo se verifica
adicionalmente mediante su lectura o validación con el parser UVL utilizado
por Flamapy. \\[0.8em]

RNF-03 &
Prueba de extremo a extremo &
\texttt{test\_step6\_generates\_one\_uvl\_model\_with\_pyodide}
(\texttt{test\_selenium.py}), que verifica que el proceso de generación se
ejecuta en el navegador mediante Pyodide y que el resultado se devuelve a la
interfaz. También está el \texttt{test\_generate\_feature\_model\_is\_discoverable\_by\_flamapy},
que verifica que Flamapy localiza esta operación mediante \texttt{DiscoverMetamodels} \\[0.8em]

RNF-04 &
Comprobación manual de usabilidad &
Recorrido manual del \textit{wizard} de configuración, comprobando que los
parámetros se presentan agrupados por funcionalidad, que las opciones
dependientes se muestran de forma coherente y que los mensajes de validación
permiten corregir una configuración inválida. La organización de las vistas
se describe en la Sección~\ref{subsec:mapa_navegacion}. \\[0.8em]

RNF-05 &
Integración continua y análisis estático &
Ejecución de los flujos de GitHub Actions configurados para el proyecto,
incluyendo las comprobaciones de calidad del código y la suite automatizada de
pruebas. Los workflows utilizados se muestran en la Figura~\ref{fig:workflows}
y la suite de pruebas se describe en la Sección~\ref{sec:pruebas_unitarias}. \\

\bottomrule

\caption{Matriz de trazabilidad entre requisitos funcionales y no funcionales y sus verificaciones.}
\label{tab:trazabilidad_rf_tests}

\end{longtable}

Un mismo caso de prueba puede aportar evidencia para varios requisitos
cuando sus comprobaciones cubren aspectos diferentes del sistema. En esos
casos, la matriz especifica en cada fila qué propiedad concreta se verifica.
Las inspecciones estructurales y las comprobaciones manuales se identifican
explícitamente y no se presentan como pruebas automatizadas.

% --- Transición al siguiente capítulo ---
\ph{Cierra el capítulo con un párrafo puente de 2--4 líneas: qué queda
establecido aquí y qué pregunta abre el capítulo siguiente. Por ejemplo:
«Con los requisitos especificados y trazados, queda fijado \emph{qué} debe
hacer el sistema. El Capítulo~\ref{chap:diseno} aborda \emph{cómo}
satisfacerlos: la arquitectura, los modelos y las decisiones de diseño que
dan forma a la solución».}

Con los requisitos funcionales, no funcionales, reglas de negocio y restricciones del proyecto definidos y trazados, queda establecido el comportamiento esperado del sistema y las condiciones que debe cumplir la solución desarrollada. El siguiente capítulo aborda cómo satisfacer estos requisitos mediante la definición de la arquitectura del sistema, los componentes principales y las decisiones de diseño adoptadas durante el desarrollo.
